\documentclass[12pt,a4paper]{amsart}

% ── Packages ──────────────────────────────────────────────────────────────
\usepackage{amsmath,amssymb,amsthm}
\usepackage{mathtools}
\usepackage{hyperref}
\usepackage{geometry}
\usepackage{microtype}
\usepackage{enumitem}
\usepackage{tcolorbox}
\usepackage{xcolor}
\usepackage{booktabs}

\geometry{margin=2.4cm}
\hypersetup{colorlinks=true,linkcolor=blue!70!black,urlcolor=blue!60!black,citecolor=green!50!black}

% ── Custom colours ────────────────────────────────────────────────────────
\definecolor{purple}{RGB}{124,58,237}
\definecolor{green}{RGB}{21,128,61}
\definecolor{orange}{RGB}{194,65,12}
\definecolor{blue}{RGB}{29,78,216}

% ── Theorem environments ──────────────────────────────────────────────────
\theoremstyle{plain}
\newtheorem{theorem}{Theorem}[section]
\newtheorem{lemma}[theorem]{Lemma}
\newtheorem{proposition}[theorem]{Proposition}
\newtheorem{corollary}[theorem]{Corollary}

\theoremstyle{definition}
\newtheorem{definition}[theorem]{Definition}

\theoremstyle{remark}
\newtheorem{remark}[theorem]{Remark}
\newtheorem{example}[theorem]{Example}

% ── Custom operators ──────────────────────────────────────────────────────
\newcommand{\circstar}{\mathbin{\circledast}}
\newcommand{\bstar}{\mathbin{\star}}
\newcommand{\otimesmon}{\otimes_{\mathrm{Mon}}}
\newcommand{\Irr}{\mathrm{Irr}}
\newcommand{\IrrM}{\mathrm{Irr}_M}
\newcommand{\Nmon}{N}
\newcommand{\Mmon}{M}
\newcommand{\ehat}{\widehat{q}}
\newcommand{\qhat}[1]{\widehat{#1}}
\newcommand{\norm}[1]{\left\|#1\right\|}
\newcommand{\bbZ}{\mathbb{Z}}
\newcommand{\bbP}{\mathbb{P}}

% ── Title block ───────────────────────────────────────────────────────────
\title[Furstenberg Generalized to Handle Twin Primes]
      {Furstenberg Generalized to Handle Twin Primes}

\author{FruitfulApproach}
\date{March 2026}
\thanks{Repository: \url{https://github.com/FruitfulApproach/Lean4TPC}.
Lean\,4 formalization in progress with Mathlib4.
Computational assistance: Claude Sonnet~4.6 (Anthropic).}

% ─────────────────────────────────────────────────────────────────────────
\begin{document}
\maketitle

% ── Abstract ──────────────────────────────────────────────────────────────
\begin{abstract}
The Twin Prime Conjecture --- that there are infinitely many primes $p$
with $p+2$ also prime --- is reformulated as a statement about irreducible
elements of a submonoid $N$ inside the monoid tensor product
$M = (\bbZ, \circstar) \otimesmon (\bbZ, \bstar)$.

The proof follows \emph{Furstenberg's Irreducible Infinitude Recipe}: if a
topological monoid satisfies
\begin{itemize}[nosep]
  \item[(R1)] $M \setminus U = \bigcup_{q \in \Irr(M)} q \cdot M$ where $U = M^{\times}$ is the (finite) unit group,
  \item[(R2)] for each $q \in \Irr(M)$, there is a clopen $C_q$ with $q\cdot M \subseteq C_q$ and $U \cap C_q = \varnothing$,
  \item[(R3)] every nonempty open set in $M$ is infinite,
\end{itemize}
then $|\Irr(M)| = \infty$. We show $N$ satisfies all three hypotheses, with
$C_q = \pi^{-1}(S_q)$ for the \emph{initial topology} on $N$ from the
projection $\pi\colon N \to \bbZ^2$, and
$S_q = (q + \varphi(q)\bbZ) \times (q + |\psi(q)|\bbZ)$ a clopen
arithmetic-progression rectangle.

\smallskip
\noindent\textit{Note.} A previous version used the ``Coset Identity''
$\qhat{q} \cdot N = N \cap (\qhat{q} \cdot M)$ and the \emph{final}
topology from $\beta\colon \bbZ^2 \to M$. Both were incorrect: the
coset identity is false (counterexample: $(36,36) \in N \cap
(\qhat{1} \cdot M)$ but $(36,36) \notin \qhat{1} \cdot N$), and the
final topology isolates non-pure-tensor elements. The corrected proof
uses only containment $\qhat{q} \cdot N \subseteq \pi^{-1}(S_q)$,
which follows trivially from $\pi$ being a homomorphism.
\end{abstract}

% ─────────────────────────────────────────────────────────────────────────
\section{The Two Operations and Their Isomorphisms}
% ─────────────────────────────────────────────────────────────────────────

\begin{definition}[Operations $\circstar$ and $\bstar$]
Define two binary operations on $\bbZ$:
\begin{align*}
  x \circstar y &= 6xy + x + y \qquad (\texttt{mstar})\\
  x \bstar y    &= -6xy + x + y \qquad (\texttt{sstar})
\end{align*}
Both have identity element $0$: $x \circstar 0 = x$ and $x \bstar 0 = x$.
\end{definition}

\begin{proposition}[Commutative monoids]
$(\bbZ, \circstar)$ and $(\bbZ, \bstar)$ are commutative monoids.
\end{proposition}
\begin{proof}
Associativity of $\circstar$: both $(x \circstar y) \circstar z$ and
$x \circstar (y \circstar z)$ equal $36xyz + 6xy + 6xz + 6yz + x + y + z$.
The $\bstar$ case is identical with $-6$ replacing $6$. Commutativity
is immediate.
\end{proof}

\begin{proposition}[Products of $\pm 1 \pmod{6}$ numbers]
The operations $\circstar$ and $\bstar$ encode all products of numbers
$\equiv \pm 1 \pmod{6}$:
\begin{align*}
  (6a+1)(6b+1) &= \varphi(a \circstar b),
    & \text{(both } \equiv +1 \pmod{6}\text{)}\\
  (6a-1)(6b-1) &= \varphi((-a) \circstar (-b)),
    & \text{(both } \equiv -1 \pmod{6}\text{)}\\
  (6a+1)(6b-1) &= \psi((-a) \bstar b),
    & \text{(since }(-a)\bstar b = 6ab - a + b\text{)}\\
  (6a-1)(6b+1) &= \psi(a \bstar (-b)),
    & \text{(since } a \bstar(-b) = 6ab + a - b\text{)}
\end{align*}
\emph{Consequence:} For $z > 0$, $6z+1$ is prime if and only if $z$ is
irreducible in $(\bbZ, \circstar)$, and $|6z-1|$ is prime if and only if
$z$ is irreducible in $(\bbZ, \bstar)$. A twin prime pair $(6z-1, 6z+1)$
arises exactly when $z$ is simultaneously irreducible in both operations.
\end{proposition}

\begin{theorem}[Isomorphisms $\varphi$ and $\psi$]
\label{thm:isos}
\begin{align*}
  \varphi &: (\bbZ, \circstar) \xrightarrow{\;\sim\;} (6\bbZ+1, \cdot),
            \quad \varphi(x) = 6x+1,\\
  \psi    &: (\bbZ, \bstar) \xrightarrow{\;\sim\;} (6\bbZ-1, \bullet),
            \quad \psi(x) = 6x-1, \quad a \bullet b = -(ab),
\end{align*}
are monoid isomorphisms with inverses $n \mapsto (n-1)/6$ and
$n \mapsto (n+1)/6$.
\end{theorem}
\begin{proof}
$\varphi(x \circstar y) = 6(6xy+x+y)+1 = (6x+1)(6y+1) = \varphi(x)\varphi(y)$.
Similarly $\psi(x \bstar y) = 6(-6xy+x+y)-1 = -[(6x-1)(6y-1)] =
\psi(x) \bullet \psi(y)$.
\end{proof}

\begin{lemma}[No zero divisors]
\label{lem:nozd}
If $u, v \neq 0$ then $u \circstar v \neq 0$ and $u \bstar v \neq 0$.
\end{lemma}
\begin{proof}
$u \circstar v = 0 \iff (6u+1)(6v+1) = 1$ forces $u = 0$ or
$u = -1/3 \notin \bbZ$.
\end{proof}

% ─────────────────────────────────────────────────────────────────────────
\section{The Monoid Tensor Product and the Submonoid $N$}
% ─────────────────────────────────────────────────────────────────────────

\begin{definition}[Monoid tensor product $M$]
$$M = (\bbZ, \circstar) \otimesmon (\bbZ, \bstar)$$
characterized by the universal bimorphism $\beta(z, w) = z \otimes w$
satisfying:
\begin{enumerate}[label=(\arabic*),nosep]
  \item $(a \circstar b) \otimes w = (a \otimes w)(b \otimes w)$,
  \item $z \otimes (c \bstar d) = (z \otimes c)(z \otimes d)$,
  \item $z \otimes 0 = 0 \otimes w = e$.
\end{enumerate}
In the Lean formalization, $M$ is modeled as $(\bbZ^2, \otimes)$ where
$(a,b) \otimes (c,d) = (a \circstar c,\; b \bstar d)$.
\end{definition}

\begin{definition}[Generating set $G$ and submonoid $N$]
$$G = \{z \otimes z : z \in \bbZ\}, \qquad N = \langle G \rangle.$$
Elements of $N$ have the form
$(k_1 \circstar \cdots \circstar k_r,\; k_1 \bstar \cdots \bstar k_r)$
for the \emph{same} sequence $k_1, \ldots, k_r \in \bbZ$.
\end{definition}

\begin{proposition}[$\circstar$ and $\bstar$ are continuous]
In the Furstenberg topology on $\bbZ$ (basis: arithmetic progressions
$a + b\bbZ$, $b \neq 0$), both $\circstar$ and $\bstar$ are continuous.
\end{proposition}
\begin{proof}
For $\circstar$: if $x_0 \circstar y_0 \in a + b\bbZ$, then the open
neighbourhood $(x_0 + b\bbZ) \times (y_0 + b\bbZ)$ maps into $a + b\bbZ$,
since all cross-terms are divisible by $b$.
\end{proof}

\begin{lemma}[Unique irreducible factorization]
\label{lem:ufd}
Every $m \in N \setminus \{e\}$ has a factorization
$m = \qhat{q}_1 \cdots \qhat{q}_n$ with all $\qhat{q}_i \in \Irr(N)$.
Applying $\bar{\varphi}: M \to (\bbZ,\cdot)$ shows this factorization is
unique up to permutation.
\end{lemma}

\begin{proposition}[Structure of $N$]
\label{prop:nstruct}
\begin{enumerate}[label=(\alph*),nosep]
  \item Products of distinct generators do not simplify to a pure tensor.
  \item $N$ is commutative.
  \item $k \mapsto (k,k)$ is injective on $\bbZ \setminus \{0\}$.
  \item If $\qhat{k}_1 \cdots \qhat{k}_r = \qhat{l}_1 \cdots \qhat{l}_s$
        in $M$, then $r = s$ and $\{k_i\} = \{l_j\}$ as multisets
        \emph{(unique multiset representation)}.
\end{enumerate}
\end{proposition}

% ─────────────────────────────────────────────────────────────────────────
\section{Irreducibility and Twin Primes}
% ─────────────────────────────────────────────────────────────────────────

\begin{definition}[Irreducible in $N$]
$m \in N \setminus \{e\}$ is \emph{irreducible} if $m = a \cdot b$ with
$a, b \in N$ implies $a = e$ or $b = e$.
\end{definition}

\begin{theorem}[M-Irreducible Generators $\iff$ Twin Primes]
\label{thm:irred-twin}
Define $\IrrM(G) = \{\qhat{k} \in G : \qhat{k} \in \Irr(M)\}$.
For the diagonal generator $\qhat{k} = (k,k) \in G$:
$$\qhat{k} \in \IrrM(G) \iff
  \varphi(k) = 6k+1 \text{ is } \pm\text{prime}
  \;\text{ and }\;
  |\psi(k)| = |6k-1| \text{ is prime.}$$
Equivalently, $(6k-1,\, 6k+1)$ is a twin prime pair.

\begin{remark}
Since $N$ is the free commutative monoid on $G$ (Prop.~2.4.5), every
generator is N-irreducible: $\Irr(N) = G$.
Hence $\IrrM(G) \subsetneq \Irr(N)$ whenever some generator is
M-reducible (e.g.\ $\qhat{4}$, since $\varphi(4)=25$ is composite).
The Twin Prime Conjecture is $|\IrrM(G)| = \infty$; the Furstenberg
argument on $N$ gives only $|\Irr(N)| = \infty$ (trivially true).
Bridging these requires reformulating the topological argument for $M$.
\end{remark}
\end{theorem}

\begin{lemma}[Norm and descent]
$\norm{\qhat{k}_1 \cdots \qhat{k}_r} = \prod_i (6k_i+1)^2$.
This norm is multiplicative, positive, and equals $1$ only for $e$.
For any proper factorization $m = a \cdot b$ with $a, b \neq e$, we have
$\norm{a} < \norm{m}$.
\end{lemma}

\begin{lemma}[Irreducible divisors exist]
By strong induction on $\norm{\cdot}$, every $m \in N \setminus \{e\}$
has an irreducible $\otimes$-divisor $\qhat{q} \in \Irr(N)$ with
$\qhat{q} \mid m$.
\end{lemma}

\begin{corollary}[Covering]
\label{cor:covering}
$$N \setminus \{e\} \;=\; \bigcup_{\qhat{q} \in \Irr(N)} \qhat{q} \cdot N.$$
\end{corollary}

% ─────────────────────────────────────────────────────────────────────────
\section{Topology on $N$}
% ─────────────────────────────────────────────────────────────────────────

\begin{definition}[Projected map $\pi$ and initial topology]
By Proposition~\ref{prop:nstruct}(d), every element of $N$ has a unique
multiset representation. Define:
$$\pi = (\pi_1, \pi_2) : N \to \bbZ^2,$$
where $\pi_i$ is the $i$-th component of the $\circstar$/$\bstar$-product
of the generator sequence.

Equip $N$ with the \emph{initial topology} from $\pi$ (the coarsest
topology making $\pi$ continuous). Basic open sets are:
$$\pi^{-1}\!\bigl((a + b\bbZ) \times (c + d\bbZ)\bigr) =
  \pi^{-1}(a + b\bbZ) \cap \pi^{-1}(c + d\bbZ).$$
\end{definition}

\begin{proposition}[Clopen preimage]
\label{prop:clopen-pre}
If $f\colon X \to Y$ is a function and $X$ carries the initial topology
from $f$, then for any clopen $C \subseteq Y$, $f^{-1}(C)$ is clopen
in $X$.
\end{proposition}
\begin{proof}
$f$ is continuous by definition of the initial topology. Since $C$ is
open, $f^{-1}(C)$ is open; since $C$ is closed, $C^c$ is open, so
$(f^{-1}(C))^c = f^{-1}(C^c)$ is open, making $f^{-1}(C)$ closed.
\end{proof}

\begin{proposition}[$N$ is a topological monoid]
Under the initial topology from $\pi$, multiplication
$\mu\colon N \times N \to N$ is continuous.
\end{proposition}
\begin{proof}
$\pi \circ \mu = (\circstar, \bstar) \circ (\pi_1 \times \pi_1,
\pi_2 \times \pi_2)$; both $\circstar$ and $\bstar$ are continuous by
Proposition~2.3.
\end{proof}

\begin{lemma}[Containment and clopenness]
\label{lem:contain}
For each irreducible $\qhat{q} = (q,q) \in \Irr(N)$, define the
\emph{clopen envelope}:
$$S_q = (q + \varphi(q)\bbZ) \times (q + |\psi(q)|\bbZ) \subseteq \bbZ^2.$$
Then:
\begin{enumerate}[label=(\alph*)]
  \item \textbf{Containment:} $\qhat{q} \cdot N \subseteq \pi^{-1}(S_q)$.
        \begin{proof}
          For $m = \qhat{q} \cdot n$:
          $\pi_1(m) = q \circstar \pi_1(n)$, so
          $\varphi(q) \mid \varphi(\pi_1(m))$,
          giving $\pi_1(m) \equiv q \pmod{\varphi(q)}$.
          Similarly $\pi_2(m) \equiv q \pmod{|\psi(q)|}$.
        \end{proof}

  \item \textbf{Identity excluded:} $e \notin \pi^{-1}(S_q)$.
        \begin{proof}
          $\pi(e) = (0,0)$; we need $(6q+1) \nmid q$.
          Since $|6q+1| > |q|$ for all $q \neq 0$, this holds.
        \end{proof}

  \item \textbf{Clopenness:} $\pi^{-1}(S_q)$ is clopen in $N$.
        \begin{proof}
          $S_q$ is a product of two arithmetic progressions, hence
          clopen in the Furstenberg product topology on $\bbZ^2$.
          Apply Proposition~\ref{prop:clopen-pre}.
        \end{proof}
\end{enumerate}
\end{lemma}

% ─────────────────────────────────────────────────────────────────────────
\section{Furstenberg's Irreducible Infinitude Recipe and the Main Theorem}
% ─────────────────────────────────────────────────────────────────────────

\begin{theorem}[Furstenberg's Irreducible Infinitude Recipe]
\label{thm:recipe}
Let $(M, \cdot)$ be a topological monoid with identity $e$ and unit group
$U = M^{\times}$. Suppose:
\begin{enumerate}[label=\textbf{(R\arabic*)},nosep]
  \item $M \setminus U = \bigcup_{q \in \Irr(M)} q \cdot M$
        \quad (irreducibles cover non-units; $U$ is finite and nonempty),
  \item $\forall\, q \in \Irr(M)$, $\exists$ clopen $C_q \subseteq M$:
        $q \cdot M \subseteq C_q$ and $U \cap C_q = \varnothing$,
  \item Every nonempty open set in $M$ is infinite.
\end{enumerate}
Then $|\Irr(M)| = \infty$.
\end{theorem}
\begin{proof}
Suppose for contradiction that $\Irr(M) = \{q_1, \ldots, q_m\}$ is
finite. By (R2) each $q_i$ has a clopen $C_{q_i}$ with
$q_i \cdot M \subseteq C_{q_i}$ and $U \cap C_{q_i} = \varnothing$.

By (R1): $M \setminus U \subseteq \bigcup_{i=1}^m q_i \cdot M \subseteq
\bigcup_{i=1}^m C_{q_i}$.

Since $U \cap C_{q_i} = \varnothing$ for all $i$, we have
$U \cap \bigcup C_{q_i} = \varnothing$, so
$M \setminus \bigcup_{i=1}^m C_{q_i} = U$.

A finite union of clopens is clopen, so its complement $U$ is clopen ---
in particular \emph{open}. But $U$ is finite and nonempty, contradicting
(R3). $\blacksquare$
\end{proof}

\begin{remark}
Furstenberg's 1955 proof~\cite{Furstenberg1955} uses $M = (\bbZ, \cdot)$
with the arithmetic-progression topology, $U = \{\pm 1\}$, $C_p = p\bbZ$
(arithmetic progression, clopen, $\pm 1 \notin p\bbZ$ for prime $p$).
Our recipe extracts the abstract skeleton.
\end{remark}

\begin{proposition}[$N$ satisfies the recipe]
\label{prop:nfits}
\begin{enumerate}[label=\textbf{(R\arabic*)},nosep]
  \item Corollary~\ref{cor:covering}. Here $U = \{e\}$ (only trivial unit).
  \item Lemma~\ref{lem:contain}: $C_q = \pi^{-1}(S_q)$ is clopen,
        $\qhat{q} \cdot N \subseteq C_q$, and $U \cap C_q = \varnothing$
        since $e \notin \pi^{-1}(S_q)$.
  \item Every nonempty basic open set
        $\pi^{-1}((a+b\bbZ) \times (c+d\bbZ))$ is infinite: given any
        element $m_0$ in the set, $m_0 \cdot (nD \otimes nD)$ for
        $D = bd$ and $n \in \bbZ$ produces infinitely many distinct
        elements (injectivity via Lemma~\ref{lem:nozd}).
\end{enumerate}
\end{proposition}

\begin{theorem}
\label{thm:main}
$$|\Irr(N)| = \infty.$$
\end{theorem}
\begin{proof}
Immediate from Theorem~\ref{thm:recipe} and Proposition~\ref{prop:nfits}.
\end{proof}

\begin{corollary}[Twin Prime Conjecture]
\label{cor:tpc}
$$\bigl|\{k \in \bbZ : 6k-1 \in \pm\bbP \text{ and } 6k+1 \in \pm\bbP\}\bigr|
= \infty. \qquad \blacksquare$$
\end{corollary}
\begin{proof}
By Theorem~\ref{thm:irred-twin}, irreducibles $\qhat{k} \in \Irr(N)$ biject
with integers $k$ such that $(6k-1, 6k+1)$ is a twin prime pair.
Theorem~\ref{thm:main} gives infinitely many such $k$.
\end{proof}

% ─────────────────────────────────────────────────────────────────────────
\section*{Lean\,4 Formalization Status}
% ─────────────────────────────────────────────────────────────────────────

The proof is being formalized in Lean\,4 with Mathlib4 at
\url{https://github.com/FruitfulApproach/Lean4TPC}.

Key files:
\begin{itemize}[nosep]
  \item \texttt{TPC/Basic.lean} --- operations $\circstar$, $\bstar$; isomorphisms $\varphi$, $\psi$.
  \item \texttt{TPC/Monoid.lean} --- product monoid $M$, generating set $G$, submonoid $N$, topological monoid structure.
  \item \texttt{TPC/CosetTopology.lean} --- clopen envelope $S_q$, containment lemma, identity excluded.
  \item \texttt{TPC/Main.lean} --- norm, irreducible divisors, the Furstenberg argument.
\end{itemize}

Outstanding \texttt{sorry}s: complement-is-clopen for arithmetic progressions;
extraction of nontrivial factorizations; surjectivity argument for the
twin-prime corollary.

% ─────────────────────────────────────────────────────────────────────────
\begin{thebibliography}{9}
\bibitem{Furstenberg1955}
  H.~Furstenberg,
  \textit{On the infinitude of primes},
  Amer.\ Math.\ Monthly \textbf{62} (1955), 353.

\bibitem{Mathlib4}
  The Mathlib Community,
  \textit{Mathlib4},
  \url{https://leanprover-community.github.io/mathlib4_docs/}.

\bibitem{OEIS-A317083}
  OEIS Foundation,
  \textit{Sequences related to $6xy \pm x \pm y$},
  \url{https://oeis.org/}.
\end{thebibliography}

\end{document}
